\documentclass[11pt]{article}
\usepackage[margin=1in]{geometry}
\usepackage{amsmath,amssymb}
\usepackage{xcolor}
\usepackage{hyperref}

\title{Parameter Justification for Carson City, NV\\SIHRS Model Case Study}
\author{SIHRS Epidemic Model}
\date{}

\begin{document}
\maketitle

\section{Justification of Parameters}
\label{JoP}

\textcolor{cyan}{For this section, I had to revert to using the symbols before renaming them. For now, I am just putting it out here the almost full picture of how the parameters were derived. Later on I guess this can be reduced to 1 or 2 paragraph of concise explanation.}

\begin{itemize}

    \item \textbf{Infection and Recovery Rates ($\beta, \gamma$)}: The basic reproduction number was set to $\mathcal{R}_0 = 1.222$, a value previously established for Carson City. To achieve this, the infection rate was set to $\beta = 0.1511$ and the base transition rate for the infectious compartment was set to $\gamma = 0.123$. While multiple combinations of $\beta$ and $\gamma$ can produce the same $\mathcal{R}_0$, these specific values were chosen to ensure stable epidemic dynamics. Rates that are too high can lead to premature fade-out, where the infected population extinguishes unrealistically early; the selected parameters produce a more sustained and realistic epidemic curve.

    \item \textbf{Hospitalization Transition Rate ($\alpha$)}: The transition rate for the hospitalized compartment was set to $\alpha = 0.111$. This value corresponds to a mean hospitalization duration of approximately 9 days ($\frac{1}{\alpha}\approx 9$), which is consistent with reported lengths of stay for COVID-19 patients.

    \item \textbf{Immunity Waning Rate ($\lambda$)}: The rate of immunity loss for the recovered compartment was $\lambda = 0.0083$. This corresponds to a mean immunity period of approximately 120 days ($ \frac{1}{\lambda} \approx 120$ days, or 4 months), reflecting the observed duration of natural immunity against COVID-19.

    \item \textbf{Hospitalization Probability ($p_{IH}$)}: The probability that an infected individual requires hospitalization was calculated as the mean ratio of hospitalized patients to active cases with a 14-day lag. Let $H(t)$ be the number of hospitalized individuals and $I_{\text{active}}(t)$ be the number of active cases on day $t$. The probability was computed over $n=72$ valid data points as:

    \[
    p_{IH} = \frac{1}{n} \sum_{k=1}^{n} \frac{H(t_k)}{I_{\text{active}}(t_k - 14)} = 0.1060
    \]

    The 14-day lag accounts for the typical period from infection to the development of severe symptoms.

    \item \textbf{Infection Fatality Probability ($p_{ID}$)}: The probability that an infected individual dies from the disease was calculated as the mean ratio of new daily deaths to active cases with a 20-day lag. Let $D_{\text{new}}(t)$ be the number of new deaths on day $t$. The formula used is:

    \[
    p_{ID} = \frac{1}{n} \sum_{k=1}^{n} \frac{D_{\text{new}}(t_k)}{I_{\text{active}}(t_k - 20)} = 0.0019
    \]

    The 20-day lag reflects the biological progression from infection to death in fatal cases.

    \item \textbf{Hospital Fatality Probability ($p_{HD}$)}: For outcomes within the hospitalized compartment, the probability of death was sourced directly from national-level data. Based on the CDC National Hospital Care Survey\cite{cdc_hospital_mortality}, the in-hospital mortality rate was set to:

    \[
    p_{HD} = 0.154
    \]

    \item \textbf{Infection to Recovery ($p_{IR}$)}: An individual in state $I$ can transition to $H$, $R$, or $D$. Therefore:

    \[
    p_{IR} = 1 - p_{IH} - p_{ID} = 1 - 0.1060 - 0.0019 = 0.8921
    \]

    \item \textbf{Hospital to Recovery ($p_{HR}$)}: An individual in state $H$ can either recover ($R$) or die ($D$). Thus:

    \[
    p_{HR} = 1 - p_{HD} = 1 - 0.154 = 0.846
    \]

    \item \textbf{Transmission Probability ($p_{SI}$)}: The probability of transmission given a contact between a susceptible and infectious individual was set to:

    \[
    p_{SI} = 1.0
    \]

    This assumption simplifies the model by embedding all transmission dynamics within the infection rate parameter, $\beta$.

    \item \textbf{Recovery to Susceptible ($p_{RS}$)}: To model the waning of natural immunity, the probability of a recovered individual returning to the susceptible state was set to:

    \[
    p_{RS} = 0.98
    \]

    This high value ensures that most individuals whose immunity wanes become fully susceptible again.

    \item \textbf{Stay Recovered ($p_{RR}$)}: An individual in state $R$ can either become susceptible ($S$) or remain recovered ($R$) upon the next clock tick.

    \[
    p_{RR} = 1 - p_{RS} = 1 - 0.98 = 0.02
    \]

\end{itemize}

\section{Summary of Parameters}

\subsection{Transition Rates}
\begin{itemize}
    \item Infection rate: $\beta = 0.1511$
    \item Infectious compartment transition rate: $\gamma = 0.123$
    \item Hospitalized compartment transition rate: $\alpha = 0.111$
    \item Immunity waning rate: $\lambda = 0.0083$
    \item Basic reproduction number: $\mathcal{R}_0 = 1.222$
\end{itemize}

\subsection{Transition Probabilities}
\begin{itemize}
    \item Susceptible to Infected: $p_{SI} = 1.0$
    \item Infected to Hospitalized: $p_{IH} = 0.1060$
    \item Infected to Death: $p_{ID} = 0.0019$
    \item Infected to Recovered: $p_{IR} = 0.8921$
    \item Hospitalized to Death: $p_{HD} = 0.154$
    \item Hospitalized to Recovered: $p_{HR} = 0.846$
    \item Recovered to Susceptible: $p_{RS} = 0.98$
    \item Recovered to Recovered: $p_{RR} = 0.02$
\end{itemize}

\subsection{Biological Timescales}
\begin{itemize}
    \item Mean infectious period: $\frac{1}{\gamma} \approx 8.1$ days
    \item Mean hospitalization duration: $\frac{1}{\alpha} \approx 9$ days
    \item Mean immunity duration: $\frac{1}{\lambda} \approx 120$ days
    \item Infection to hospitalization lag: 14 days
    \item Infection to death lag: 20 days
\end{itemize}

\section{Data Sources}

\begin{itemize}
    \item \textbf{Hospitalization Data}: HealthData.gov COVID-19 Reported Patient Impact and Hospital Capacity by Facility (August 2020 - December 2021)
    \item \textbf{Case and Death Data}: The New York Times COVID-19 Data Repository (March 2020 - December 2021)
    \item \textbf{Hospital Mortality Rate}: CDC National Hospital Care Survey
    \item \textbf{Basic Reproduction Number}: Previously established value for Carson City, NV
\end{itemize}

\section{Notes}

All empirically-derived parameters ($p_{IH}$, $p_{ID}$) are specific to Carson City, Nevada, and reflect the local COVID-19 pandemic dynamics during the March 2020 - December 2021 period. These values should be adjusted for other geographic regions or time periods based on local data.

\end{document}

